% !PRINTED 728
\weberpage{728}
%<<CONT
übrigen aber in umgekehrter Ordnung und mit verändertem Vorzeichen schreiben.
Dabei sind $\frac{1}{2}\nu'$ Vorzeichenänderungen nöthig,
und wir erhalten daher für das Product (12) die durch folgende
Gleichung definirte Determinante, deren absoluten Werth wir
mit $T$ bezeichnen:
\weq{14}{(-1)^{\frac{1}{2}\nu'}
\begin{vmatrix}
\lambda_0, & \lambda_1, & \ldots, & \lambda_{\nu'-2}, & \lambda_{\nu'-1}\\
\lambda_1, & \lambda_2, & \ldots, & \lambda_{\nu'-1}, & -\lambda_0\\
\lambda_2, & \lambda_3, & \ldots, & -\lambda_0, & -\lambda_1\\
\cdots & \cdots & \cdots & \cdots & \cdots\\
\lambda_{\nu'-1}, & -\lambda_0, & \ldots, & -\lambda_{\nu'-3}, & -\lambda_{\nu'-2}
\end{vmatrix} = \pm T,}
worin rechts von der von links unten nach rechts oben laufenden
Diagonalreihe die negativen Zeichen stehen.

Die Determinante $T$ ist hier so geordnet, dass jede Zeile
aus der vorangehenden durch die Substitution $(r, r^{5^\beta})$ entsteht.
Jede Colonne enthält also ein System conjugirter Logarithmen.

Nach (12) ist jetzt
\weq{15}{B = \frac{T}{D}.}
Um über die Natur dieses Resultates Klarheit zu bekommen,
ist es nöthig, die Zahl $D$ etwas genauer zu betrachten, die durch
\S.\,195, (12), \S.\,197, (6) definirt war; und dies erfordert ein Eingehen
 auf die Theorie der Einheiten des Körpers $H_m$.

\begin{center}
\S.\,200.\\
\gesperrt{Normaleinheiten in $H_m$.}
\end{center}

Wir führen jetzt eine vereinfachende Bezeichnung ein, indem wir
\weq{1}{r^{5^\beta} = r_\beta}
setzen, und $\beta$ ein volles Restsystem nach dem Modul $\nu$ durchlaufen
 lassen. Es folgt hieraus
\weq{2}{r_{\beta+\nu'} = -r_\beta,}
und in der Form $(r, r_\beta)$ sind dann zwar nicht alle Substitutionen
der Gruppe des Körpers $\Omega_m$ enthalten, sondern es müssen noch
die Substitutionen $(r, r_\beta^{-1})$ hinzukommen; wohl aber liefert uns
$(r, r_\beta)$ alle Substitutionen der Gruppe von $H_m$.

Wir bezeichnen jetzt mit $\mathfrak{E}(r)$ eine Einheit des Körpers $H_m$.
Unter diesen sind auch die Einheiten des Körpers $H_\mu$ (als
%CONT>>
